\section{Пустота КС-языка}
	\tikzsetfigurename{CFG_Emptiness_}

\begin{theorem}
    Существует алгоритм, определяющий, является ли язык, порождаемый КС грамматикой, пустым.
\end{theorem}

\begin{proof}
    Следующая лемма утверждает, что если в КС языке есть выводимое слово, то существует другое выводимое слово с деревом вывода не глубже количества нетерминалов грамматики.
    Для доказательства теоремы достаточно привести алгоритм, последовательно строящий все деревья глубины не больше количества нетерминалов грамматики, и проверяющий, являются ли такие деревья деревьями вывода.
    Если в результате работы алгоритма не удалось построить ни одного дерева, то грамматика порождает пустой язык.
\end{proof}

\begin{lemma}
    Если в данной грамматике выводится некоторая цепочка, то существует цепочка, дерево вывода которой не содержит ветвей длиннее $m$, где $m$~--- количество нетерминалов грамматики.
\end{lemma}

\begin{proof}
    Рассмотрим дерево вывода цепочки $\omega$. Если в нем есть 2 узла, соответствующих одному нетерминалу A, обозначим их $n_1$ и $n_2$.

    Предположим, $n_1$ расположен ближе к корню дерева, чем $n_2$.

    Вывод цепочки $\omega$ имеет следующий вид:
    \[S \derives \alpha A_{n_1} \beta \derives \alpha \omega_1 \beta; S \derives \alpha \gamma A_{n_2} \delta \beta \derives \alpha \gamma \omega_2 \delta \beta \derives \omega,\]
    при этом $\omega_2$ является подцепочкой $\omega_1$.

    Заменим в изначальном дереве узел $n_1$ на $n_2$. Полученное дерево является деревом вывода цепочки $\alpha \omega_2 \delta$.

    Повторяем процесс замены одинаковых нетерминалов до тех пор, пока в дереве не останутся только уникальные нетерминалы.

    В итоге, в полученном дереве не может быть ветвей длины большей, чем $m$, и по построению оно является деревом вывода.
\end{proof}
